% SGA 3, Expose II, Loop 1 English reconstruction.
% Scope: Section 4.1 through the end of Example 4.4.2, stopping immediately
% before Proposition 4.5.
% Authority: Polo--Gille Exp2-14oct24.pdf, component-local pages 22--34
% (combined-reader pages 80--92; printed pages 70--82). Force-OCR is a
% drafting locator; English wording was compared with Jacob Reinhold's
% CC BY 4.0 Markdown lineage.

\setcounter{footnote}{48}

% Source page: Exp2 component-local p. 22; combined-reader p. 80; printed p. 70.
\section*{\SGAInlineRefTarget{sga3:ii:section:4}4. 群的切空间---Lie 代数}
\addcontentsline{toc}{section}{4. 群的切空间---Lie 代数}

\subsection*{\SGAInlineRefTarget{sga3:ii:subsection:4-1}4.1.}
设 \(G\) 是 \(S\) 上的群函子. 由 \SGARefLink{sga3:ii:corollary:3-9-0-2}{3.9.0.2},
\(T_{G/S}(\mathcal M)\) 与
\(\underline{\operatorname{Lie}}(G/S,\mathcal M)\) 均带有 \(S\) 上的群结构,
且有群态射

\SGAThreeNativeDiagram[0.52\textwidth]
  {figures/exp2/exp2_p022_split_tangent_sequence.tex}
  {群函子 \(G\) 的分裂切序列.}
  {fig:exp2-split-tangent-sequence}
  {Polo--Gille Exp2-14oct24.pdf, 组件内第 22 页; 合卷第 80 页; 印刷页码 70}

由定义, \(i\) 将 \(\underline{\operatorname{Lie}}(G/S,\mathcal M)\) 等同于
\(p\) 的核, 而 \(s\) 是 \(p\) 的截面.
由 \SGARefLink{sga3:i:proposition:2-3-7}{I~2.3.7} 可知, 这一序列将
\(T_{G/S}(\mathcal M)\) 等同于
\(\underline{\operatorname{Lie}}(G/S,\mathcal M)\) 与 \(G\) 的半直积.

\SGAInlineRefTarget{sga3:ii:definition:4-1-a}\subsection*{\SGARefLink{sga3:ii:definition:4-1-a}{定义 4.1.A}.}
相应的 \(G\) 在 \(\underline{\operatorname{Lie}}(G/S,\mathcal M)\) 上的作用记为
\[
  \operatorname{Ad}:G\longrightarrow
  \underline{\operatorname{Aut}}_{S\text{-gr.}}
    \bigl(\underline{\operatorname{Lie}}(G/S,\mathcal M)\bigr)
\]
并称为 \(G\) 相对于 \(\mathcal M\) 的\emph{伴随表示 (adjoint representation)}.
于是, 由定义, 对 \(x\in G(S')\) 和
\(X\in\underline{\operatorname{Lie}}(G/S,\mathcal M)(S')\), 有
\[
  \operatorname{Ad}(x)X
  =i^{-1}\bigl(s(x)i(X)s(x)^{-1}\bigr).
\]
若 \(G\) 交换, 则 \(T_{G/S}(\mathcal M)\) 也交换, 且
\(\operatorname{Ad}(x)X=X\).
\footnote{编者注: 增添了编号 \SGARefLink{sga3:ii:definition:4-1-a}{4.1.A}
和 \SGARefLink{sga3:ii:definition:4-1-b}{4.1.B}, 以突出这两个定义.}

\SGAInlineRefTarget{sga3:ii:definition:4-1-b}\subsection*{\SGARefLink{sga3:ii:definition:4-1-b}{定义 4.1.B}.}
设 \(G,H\) 是 \(S\) 上的群函子, \(f:G\to H\) 是群态射.
函子性给出一个与典范截面相容的正合序列态射.

\SGAThreeNativeDiagram[0.66\textwidth]
  {figures/exp2/exp2_p022_derived_exact_sequences.tex}
  {分裂切序列的函子性.}
  {fig:exp2-derived-exact-sequences}
  {Polo--Gille Exp2-14oct24.pdf, 组件内第 22 页; 合卷第 80 页; 印刷页码 70}

态射 \(L(f)\), 也记为 \(\underline{\operatorname{Lie}}(f)\),
称为 \(f\) 的\emph{导出态射 (derived morphism)}.
\footnote{编者注: 若 \(f\) 为单态射, 则 \(L(f)\) 和 \(T(f)\) 也为单态射;
参见 \SGARefLink{sga3:ii:proposition:3-7}{3.7}.}

\SGAInlineRefTarget{sga3:ii:remark:4-1-c}\subsection*{\SGARefLink{sga3:ii:remark:4-1-c}{注 4.1.C}.}
若 \(G/S\) 和 \(H/S\) 满足条件 \(E\), 则 \(L(f)\) 保持由
“关于 \(\mathcal M\) 的函子性”得到的 \(\mathcal O_S\)-模结构
(参见 \SGARefLink{sga3:ii:proposition:3-6}{3.6}).
\footnote{编者注: 此注为增补内容.}

\SGAInlineRefTarget{sga3:ii:proposition:4-1-1}\subsection*{\SGARefLink{sga3:ii:proposition:4-1-1}{命题 4.1.1}.}
设 \(g\in G(S)\). 则
\[
  \operatorname{Ad}(g):
  \underline{\operatorname{Lie}}(G/S,\mathcal M)
  \longrightarrow
  \underline{\operatorname{Lie}}(G/S,\mathcal M)
\]
是 \(\operatorname{Int}(g):G\to G\) 的导出态射.

事实上, \(\operatorname{Ad}(g)X=i^{-1}(\operatorname{Int}(g)i(X))\),
根据导出态射的定义, 这正是 \(L(\operatorname{Int}(g))(X)\).

设 \(G/S\) 满足 \(E\). 由\SGARefLink{sga3:ii:proposition:3-9}{命题~3.9},
刚刚在 \(\underline{\operatorname{Lie}}(G/S,\mathcal M)\) 上定义的群结构,
恰好是由其 \(\mathcal O_S\)-模结构 (由 (E) 定义) 诱导的群结构.
因此, \SGARefLink{sga3:ii:proposition:4-1-1}{命题~4.1.1} 与
\(\mathcal O_S\)-作用的函子性 (参见 \SGARefLink{sga3:ii:proposition:3-6}{3.6})
给出下述推论.

% Source page: Exp2 component-local p. 23; combined-reader p. 81; printed p. 71.
\SGAInlineRefTarget{sga3:ii:corollary:4-1-1-1}\subsection*{\SGARefLink{sga3:ii:corollary:4-1-1-1}{推论 4.1.1.1}.}
设 \(G/S\) 满足 \(E\). 则 \(\operatorname{Ad}\) 将 \(G\) 映入子群
\[
  \underline{\operatorname{Aut}}_{\mathcal O_S\text{-mod.}}
    \bigl(\underline{\operatorname{Lie}}(G/S,\mathcal M)\bigr)
\]
它是 \(\underline{\operatorname{Aut}}_{S\text{-gr.}}
(\underline{\operatorname{Lie}}(G/S,\mathcal M))\) 的子群. 换言之, 对每个 \(g\in G(S')\), 自同构 \(\operatorname{Ad}(g)\) 保持
\(\underline{\operatorname{Lie}}(G'/S',\mathcal M)\) 上的
\(\mathcal O_{S'}\)-模结构, 其中 \(G'=G\times_SS'\).
因此, \(\operatorname{Ad}\) 是 \(G\) 在 \(\mathcal O_S\)-模
\(\underline{\operatorname{Lie}}(G/S,\mathcal M)\) 上的一个
\emph{线性表示} (参见 \SGARefLink{sga3:i:definition:3-2}{I~3.2}).

\SGAInlineRefTarget{sga3:ii:remark:4-1-1-2}\subsection*{\SGARefLink{sga3:ii:remark:4-1-1-2}{注 4.1.1.2}.}
\textit{a)} 函子 \(G/S\) 满足 \(E\), 当且仅当对于每一对有限自由
\(\mathcal O_S\)-模 \((\mathcal M,\mathcal N)\), 下图是笛卡尔图;
该图是将 \(\underline{\operatorname{Lie}}(G/S,-)\) 用于
\SGARefLink{sga3:ii:definition:2-1}{2.1} 的图 \((*)\) 得到的.

\SGAThreeNativeDiagram[0.45\textwidth]
  {figures/exp2/exp2_p023_condition_e_cartesian.tex}
  {条件 \(E\) 的笛卡尔判据.}
  {fig:exp2-condition-e-cartesian}
  {Polo--Gille Exp2-14oct24.pdf, 组件内第 23 页; 合卷第 81 页; 印刷页码 71}

\textit{b)} 设 \(G/S\) 满足 \(E\). 群律 \(\pi:G\times_SG\to G\) 的导出态射,
就是 \(\underline{\operatorname{Lie}}(G/S,\mathcal M)\) 上的加法律.
注意, \(\pi\) 并非群态射, 但 \(\pi(e,e)=e\); 因而导出态射 \(L(\pi)\) 将
\[
  T^{(e,e)}_{(G\times_SG)/S}(\mathcal M)
  =\underline{\operatorname{Lie}}(G/S,\mathcal M)
   \times_S\underline{\operatorname{Lie}}(G/S,\mathcal M)
\]
映入 \(\underline{\operatorname{Lie}}(G/S,\mathcal M)\);
参见 \SGARefLink{sga3:ii:proposition:3-7}{3.7} 和
\SGARefLink{sga3:ii:proposition:3-8}{3.8}. 同样, 对 \(n\in\mathbb Z\),
以 \(n_G:G\to G\) 表示 \(S\)-函子态射 \(g\mapsto g^n\).
则 \(L(n_G)\) 是 \(\underline{\operatorname{Lie}}(G/S)\) 上的乘以 \(n\) 的映射;
参见 \SGARefLink{sga3:ii:corollary:3-9-4}{3.9.4}.

\subsection*{\SGAInlineRefTarget{sga3:ii:subsection:4-1-2-0}4.1.2.0.}
考虑 \(S\)-函子
\(
  \underline{\operatorname{Hom}}_{G/S}
    (G,T_{G/S}(\mathcal M)).
\)\footnote{编者注: 下文对原文作了扩充, 以明确展示 \(G\) 在函子
\(\underline{\operatorname{Hom}}_{G/S}(G,T_{G/S}(\mathcal M))\) 及
\(\underline{\operatorname{Hom}}_S
(G,\underline{\operatorname{Lie}}(G/S,\mathcal M))\) 上的作用.}
对每个 \(S'\to S\), 有
\(T_{G/S}(\mathcal M)_{S'}\simeq T_{G_{S'}/S'}(\mathcal M)\)
(参见 \SGARefLink{sga3:ii:proposition:3-4}{3.4}), 因而
\[
\begin{aligned}
 &\underline{\operatorname{Hom}}_{G/S}
   (G,T_{G/S}(\mathcal M))(S')\\
 &\qquad\simeq
 \operatorname{Hom}_{G_{S'}}
   \bigl(G_{S'},T_{G_{S'}/S'}(\mathcal M)\bigr)
 =\Gamma\bigl(T_{G_{S'}/S'}(\mathcal M)/G_{S'}\bigr).
\end{aligned}
\]
先注意以下关于 \(S'\) 函子性的同构:
\begin{equation}
 \operatorname{Hom}_{S'}
  \bigl(G_{S'},\underline{\operatorname{Lie}}(G_{S'}/S',\mathcal M)\bigr)
 \xrightarrow{\sim}
 \Gamma\bigl(T_{G_{S'}/S'}(\mathcal M)/G_{S'}\bigr),
 \tag{*}\label{sga3:ii:subsection:4-1-2-0:equation:star}
\end{equation}
它把 \(f:G_{S'}\to
\underline{\operatorname{Lie}}(G_{S'}/S',\mathcal M)\) 对应到截面
\(s_f:G_{S'}\to T_{G_{S'}/S'}(\mathcal M)\), 后者对每个
\(S''\to S'\) 和 \(g\in G(S'')\) 由下式定义:
\[
  s_f(g)=i(f(g))s(g).
\]

% Source page: Exp2 component-local p. 24; combined-reader p. 82; printed p. 72.
设 \(h\) 是 \(S'\) 上函子 \(G_{S'}\) 的一个自同构, 不要求它保持群结构.
对 \(T_{G_{S'}/S'}(\mathcal M)\) 的每个截面 \(\tau\), 通过结构迁移得到一个截面
\(h(\tau)\): 它是使下图交换的唯一截面.

\SGAThreeNativeDiagram[0.43\textwidth]
  {figures/exp2/exp2_p024_transport_square.tex}
  {自同构 \(h\) 对切截面的转移.}
  {fig:exp2-transport-square}
  {Polo--Gille Exp2-14oct24.pdf, 组件内第 24 页; 合卷第 82 页; 印刷页码 72}

特别地, 取 \(h=t_x\) 为元素 \(x\in G(S')\) 给出的右平移;
于是对所有 \(g\in G(S'')\), \(S''\to S'\), 有 \(h(g)=t_x(g)=g\cdot x\).
此时有
\[
  t_x(s_f)=s_{t_x(f)},
\]
其中 \(t_x(f):G_{S'}\to
\underline{\operatorname{Lie}}(G_{S'}/S',\mathcal M)\) 由下式给出:
\[
  t_x(f)(g)=f(g\cdot x^{-1}).
\]

让 \(G\) 通过右平移按如下方式作用于
\[
 \underline{\operatorname{Hom}}_{G/S}(G,T_{G/S}(\mathcal M))
 \quad\text{和}\quad
 \underline{\operatorname{Hom}}_S
   (G,\underline{\operatorname{Lie}}(G/S,\mathcal M))
\]
对 \(S'\to S\) 及 \(x\in G(S')\), 设
\[
\begin{aligned}
 \sigma&\in\Gamma(T_{G_{S'}/S'}(\mathcal M)/G_{S'}),\\
 f&\in\operatorname{Hom}_{S'}
 (G_{S'},\underline{\operatorname{Lie}}(G_{S'}/S',\mathcal M)).
\end{aligned}
\]
对每个 \(g\in G(S'')\), \(S''\to S'\), 置
\[
 (\sigma\cdot x)(g)=\sigma(g\cdot x^{-1})\cdot s(x),
 \qquad
 (f\cdot x)(g)=f(g\cdot x^{-1})
\]
同构 \((*)\) 保持这些 \(G\)-作用.

因此, 在这一同构下,
\[
 \underline{\operatorname{Hom}}_{G/S}
   (G,T_{G/S}(\mathcal M))^G(S')
\]
的元素对应于从 \(G_{S'}\) 到
\(\underline{\operatorname{Lie}}(G_{S'}/S',\mathcal M)\) 的常态射,
即经过 \(G_{S'}\to S'\) 分解的态射; 等价地,
它们对应于 \(\underline{\operatorname{Lie}}(G/S,\mathcal M)(S')\) 的元素.

\paragraph{术语.}
\(\underline{\operatorname{Hom}}_{G/S}
(G,T_{G/S}(\mathcal M))^G(S')\) 的元素称为
\(T_{G_{S'}/S'}(\mathcal M)\) 的右平移不变截面.

\SGAInlineRefTarget{sga3:ii:proposition:4-1-2}\subsection*{\SGARefLink{sga3:ii:proposition:4-1-2}{命题 4.1.2}.}
映射
\[
 \underline{\operatorname{Lie}}(G/S,\mathcal M)(S)
 \longrightarrow \Gamma(T_{G/S}(\mathcal M)/G),
 \qquad X\longmapsto (x\mapsto X\cdot x),
\]
是到右平移不变截面集合的双射.
\footnote{只要注意到 \(G\) 的右平移不变自同构就是左平移,
便可更简单地得到陈述 4.1.2, 4.1.3 和 4.1.4.
编者注: 例如参见 \textup{[DG70]}, \S II.4 中命题 4.6 与 2.5 的证明.}

% Source page: Exp2 component-local p. 25; combined-reader p. 83; printed p. 73.
同样, 让 \(G\) 按下式在
\(\underline{\operatorname{End}}_{I_S(\mathcal M)/S}
(G_{I_S(\mathcal M)})\) 上右作用:
\[
 (u\cdot x)(g)=u(g\cdot x^{-1})\cdot x,
\]
其中 \(S'\to S\), \(x\in G(S')\),
\(u\in\operatorname{End}_{I_{S'}(\mathcal M)}
(G_{I_{S'}(\mathcal M)})\), 且 \(g\in G(S'')\),
\(S''\to I_{S'}(\mathcal M)\). 3.12.1 中的态射
\[
 \underline{\operatorname{Hom}}_{G/S}(G,T_{G/S}(\mathcal M))
 \longrightarrow
 \underline{\operatorname{End}}_{I_S(\mathcal M)/S}
   (G_{I_S(\mathcal M)})
\]
保持 \(G\)-作用. 因此, 对每个 \(S'\to S\), 它在
\(\Gamma(T_{G_{S'}/S'}(\mathcal M)/G_{S'})\) 与下述自同态的集合之间诱导出一个双射:
\(G_{I_{S'}(\mathcal M)}\) 的 \(I_{S'}(\mathcal M)\)-自同态 \(u\),
它在 \(G\) 上诱导恒等映射, 且“与右平移可交换”, 其含义是对每个 \(S''\to S'\)
及 \(x\in G(S'')\), 有 \(u_{S''}\cdot x=u_{S''}\).

\SGAInlineRefTarget{sga3:ii:proposition:4-1-3}\subsection*{\SGARefLink{sga3:ii:proposition:4-1-3}{命题 4.1.3}.}
在 \(\underline{\operatorname{Lie}}(G/S,\mathcal M)(S)\) 与下述自同态的集合之间,
存在一个关于 \(G\) 函子性的双射:
\(G_{I_S(\mathcal M)}\) 的 \(I_S(\mathcal M)\)-自同态,
它在 \(G\) 上诱导恒等映射, 并在刚才说明的意义下与 \(G\) 的右平移可交换.

结合 \SGARefLink{sga3:ii:proposition:3-13}{3.13}, 得到下述定理.

\SGAInlineRefTarget{sga3:ii:theorem:4-1-4}\subsection*{\SGARefLink{sga3:ii:theorem:4-1-4}{定理 4.1.4}.}
设 \(G\) 是取值于群的 \(S\)-函子, 且 \(G/S\) 满足 \(E\).
则 \(\underline{\operatorname{Lie}}(G/S,\mathcal M)(S)\)
与下述自同构组成的群关于 \(G\) 函子性地等同:
\(G_{I_S(\mathcal M)}\) 的 \(I_S(\mathcal M)\)-自同构,
它在 \(G\) 上诱导恒等映射, 并与 \(G\) 的右平移可交换.

当 \(\mathcal M=\mathcal O_S\) 时, 这就得到群的 Lie 代数的一个经典定义.

\subsection*{\SGAInlineRefTarget{sga3:ii:subsection:4-2-0}4.2.0.}
在继续讨论之前, 先证明 \SGARefLink{sga3:ii:proposition:3-11}{3.11} 的另外一些推论.
\footnote{编者注: \SGARefLink{sga3:ii:subsection:4-2-0}{4.2.0 段}是增补内容;
其结果在原文的 \SGARefLink{sga3:ii:subsection:4-2}{4.2} 和
\SGARefLink{sga3:ii:theorem:4-7}{4.7} 中被隐含地使用.}
如\SGARefLink{sga3:ii:section:1}{第~1~节}, 设 \(X,Y\) 在 \(Z\) 上,
而 \(Z\) 在 \(S\) 上. 正如 \SGARefLink{sga3:ii:proposition:3-11}{3.11} 所说明的,
同构 \SGARefLink{fig:exp2-hom-adjunction-triangle}{1.1(4)}

\SGAThreeNativeDiagram[0.78\textwidth]
  {figures/exp2/exp2_p025_hom_adjunction_triangle.tex}
  {\SGARefLink{fig:exp2-hom-adjunction-triangle}{1.1(4)} 中的伴随同构.}
  {fig:exp2-hom-adjunction-p25}
  {Polo--Gille Exp2-14oct24.pdf, 组件内第 25 页; 合卷第 83 页; 印刷页码 73}

诱导下述同构 \(\theta\).

\SGAThreeNativeDiagram[0.78\textwidth]
  {figures/exp2/exp2_p025_tangent_adjunction_triangle.tex}
  {所诱导的切空间同构 \(\theta\).}
  {fig:exp2-tangent-adjunction-p25}
  {Polo--Gille Exp2-14oct24.pdf, 组件内第 25 页; 合卷第 83 页; 印刷页码 73}

% Source page: Exp2 component-local p. 26; combined-reader p. 84; printed p. 74.
由 \SGARefLink{sga3:ii:remark:1-3}{1.3}, 若 \(Y\) 是 \(Z\)-群,
则对每个 \(V\to Z\), \(\underline{\operatorname{Hom}}_Z(V,Y)\) 也是 \(Z\)-群;
特别地, 可以取 \(V=I_Z(\mathcal M)\). 明确地说, 若 \(Z''\to Z'\to Z\)
且 \(\varphi,\psi\in\operatorname{Hom}_Z(V_{Z'},Y)(Z')\), 则它们的积由下式定义:
\[
  (\varphi\cdot\psi)(v)=\varphi(v)\psi(v)
  \qquad (v\in V_{Z'}(Z'')).
\]

现设 \(X\) 和 \(Y\) 都是 \(Z\)-群.

\SGAInlineRefTarget{sga3:ii:definition:4-2-0-1}\subsection*{\SGARefLink{sga3:ii:definition:4-2-0-1}{定义 4.2.0.1}.}
以 \(\underline{\operatorname{Hom}}_{(Z/S)\text{-gr.}}(X,Y)\) 表示
\(\underline{\operatorname{Hom}}_{Z/S}(X,Y)\) 的下述子函子:
对每个 \(S'\to S\),
\begin{equation}
 \underline{\operatorname{Hom}}_{(Z/S)\text{-gr.}}(X,Y)(S')
 =\operatorname{Hom}_{Z_{S'}\text{-gr.}}(X_{S'},Y_{S'}).
 \tag{3}\label{sga3:ii:subsection:4-2-0:equation:3}
\end{equation}
将 \(Y\) 换成 \(Z\)-群 \(T_{Y/Z}(\mathcal M)\) 时, 也采用同一定义.

在上面的同构 (2) 下,
\(T_{\underline{\operatorname{Hom}}_{(Z/S)\text{-gr.}}(X,Y)/S}
(\mathcal M)(S')\) 对应于关于 \(X\) 为乘法性的 \(Z_{S'}\)-态射
\[
 \varphi:X_{S'}\times_{S'}I_{S'}(\mathcal M)\longrightarrow Y_{S'}
\]
即
\[
 \varphi(x_1x_2,m)=\varphi(x_1,m)\varphi(x_2,m).
\]
这些态射又对应于 \(Z_{S'}\)-群态射
\(X_{S'}\to T_{Y/Z}(\mathcal M)_{S'}\). 因此:

\SGAInlineRefTarget{sga3:ii:proposition:4-2-0-2}\subsection*{\SGARefLink{sga3:ii:proposition:4-2-0-2}{命题 4.2.0.2}.}
设 \(X,Y\) 是 \(Z\)-群, 其中 \(Z\) 在 \(S\) 上.
存在关于 \(\mathcal M\) 函子性的 \(S\)-函子同构
\[
 T_{\underline{\operatorname{Hom}}_{(Z/S)\text{-gr.}}(X,Y)/S}
   (\mathcal M)
 \xrightarrow{\sim}
 \underline{\operatorname{Hom}}_{(Z/S)\text{-gr.}}
   (X,T_{Y/Z}(\mathcal M)).
\]

取 \(Z=S\), 得到如下推论. 若 \(Y\) 是交换 \(S\)-群,
则 \(T_{Y/S}(\mathcal M)\) 也如此, 因而
\(H=\underline{\operatorname{Hom}}_{S\text{-gr.}}(X,Y)\),
\(\underline{\operatorname{Hom}}_{S\text{-gr.}}
(X,T_{Y/S}(\mathcal M))\) 及 \(T_{H/S}(\mathcal M)\) 也都如此.

\SGAInlineRefTarget{sga3:ii:corollary:4-2-0-3}\subsection*{\SGARefLink{sga3:ii:corollary:4-2-0-3}{推论 4.2.0.3}.}
设 \(X,Y\) 是 \(S\)-群. 存在关于 \(\mathcal M\) 函子性的 \(S\)-函子同构
\[
 T_{\underline{\operatorname{Hom}}_{S\text{-gr.}}(X,Y)/S}
   (\mathcal M)
 \xrightarrow{\sim}
 \underline{\operatorname{Hom}}_{S\text{-gr.}}
   (X,T_{Y/S}(\mathcal M)).
\]
若 \(Y\) 交换, 则它们还是交换 \(S\)-群的同构.

\SGAInlineRefTarget{sga3:ii:definition:4-2-0-4}\subsection*{\SGARefLink{sga3:ii:definition:4-2-0-4}{定义 4.2.0.4}.}
若 \(Y\) 是 \(\mathcal O_S\)-模, 则函子 \(T_{Y/S}(\mathcal M)\),
相应地 \(\underline{\operatorname{Lie}}(Y/S,\mathcal M)\),
带有一个由 \(Y\) 的结构诱导的 \(\mathcal O_S\)-模结构.
配备这一结构后, 将它记为 \(T'_{Y/S}(\mathcal M)\),
相应地记为 \(\underline{\operatorname{Lie}}'(Y/S,\mathcal M)\).
\footnote{编者注: 撇号用来区分由 \(Y\) 的结构诱导的 \(\mathcal O_S\)-模结构,
与将要从条件 \(E\) 得到的模结构.}

因此, 若 \(X,Y\) 是 \(\mathcal O_S\)-模, 则
\(T'_{Y/S}(\mathcal M)=
\underline{\operatorname{Hom}}_S(I_S(\mathcal M),Y)\).
若 \(H=\underline{\operatorname{Hom}}_{\mathcal O_S\text{-mod.}}(X,Y)\), 则
\(\underline{\operatorname{Hom}}_{\mathcal O_S\text{-mod.}}
(X,T'_{Y/S}(\mathcal M))\) 和 \(T'_{H/S}(\mathcal M)\)
也带有 \(\mathcal O_S\)-模结构.

% Source page: Exp2 component-local p. 27; combined-reader p. 85; printed p. 75.
\SGAInlineRefTarget{sga3:ii:corollary:4-2-0-5}\subsection*{\SGARefLink{sga3:ii:corollary:4-2-0-5}{推论 4.2.0.5}.}
若 \(X,Y\) 是 \(\mathcal O_S\)-模, 则存在关于 \(\mathcal M\) 函子性的
\(\mathcal O_S\)-模同构
\[
 T'_{\underline{\operatorname{Hom}}_{\mathcal O_S\text{-mod.}}(X,Y)/S}
   (\mathcal M)
 \xrightarrow{\sim}
 \underline{\operatorname{Hom}}_{\mathcal O_S\text{-mod.}}
   (X,T'_{Y/S}(\mathcal M)).
\]

\SGAInlineRefTarget{sga3:ii:definition:4-2-a}\subsection*{\SGARefLink{sga3:ii:definition:4-2-a}{定义 4.2.A}.}
设 \(X,L\) 是 \(S\)-群, \(X\) 通过群自同构作用于 \(L\)
(参见 \SGARefLink{sga3:i:definition:2-3-5}{I~2.3.5}).
以下式定义 \(\underline{\operatorname{Hom}}_S(X,L)\) 的子函子
\(\underline Z^1_S(X,L)\):
\[
\begin{aligned}
 \underline Z^1_S(X,L)(S')
 =\{\,\varphi\in\operatorname{Hom}_{S'}(X_{S'},L_{S'})\mid{}
 &\varphi(x_1x_2)
   =\varphi(x_1)\bigl(x_1\cdot\varphi(x_2)\bigr),\\[-2pt]
 &x_1,x_2\in X(S''),\ S''\to S'\,\}.
\end{aligned}
\]
它称为从 \(X\) 到 \(L\) 的\emph{交叉同态函子 (functor of crossed homomorphisms)}.
\footnote{编者注: 下文对原文作了扩充; 特别是增添了定义 4.2.A 与注 4.2.B.}

\SGAInlineRefTarget{sga3:ii:remark:4-2-b}\subsection*{\SGARefLink{sga3:ii:remark:4-2-b}{注 4.2.B}.}
\textit{a)} 若 \(L'\) 是另一个 \(S\)-群, 且 \(X\) 也通过群自同构作用于它,
则
\[
 \underline Z^1_S(X,L\times_SL')
 \simeq \underline Z^1_S(X,L)\times_S\underline Z^1_S(X,L').
\]

\textit{b)} 若 \(L\) 是 \(X\)-\(\mathcal O_S\)-模,
则 \(\underline Z^1_S(X,L)\) 是 \SGARefLink{sga3:i:subsection:5-1}{I~5.1} 中定义的微分
\[
 \partial:\underline{\operatorname{Hom}}_S(X,L)
 \longrightarrow\underline{\operatorname{Hom}}_S(X^2,L)
\]
的核. 特别地, 此时 \(\underline Z^1_S(X,L)\) 是 \(\mathcal O_S\)-模.

\subsection*{\SGAInlineRefTarget{sga3:ii:subsection:4-2}4.2.}
设 \(u:X\to Y\) 是 \(S\)-群态射. 我们要描述
\(L^u_{\underline{\operatorname{Hom}}_{S\text{-gr.}}(X,Y)/S}
(\mathcal M)\). 首先, \SGARefLink{sga3:ii:corollary:3-11-3}{3.11.3}
给出关于 \(\mathcal M\) 函子性的 \(S\)-函子同构
\begin{equation}
 L^u_{\underline{\operatorname{Hom}}_S(X,Y)/S}(\mathcal M)
 \xrightarrow{\sim}
 \underline{\operatorname{Hom}}_{Y/S}(X,T_{Y/S}(\mathcal M)).
 \tag{\(\dagger\)}\label{sga3:ii:subsection:4-2:equation:dagger}
\end{equation}
由于 \(Y\) 是 \(S\)-群,
\[
 T_{Y/S}(\mathcal M)
 =\underline{\operatorname{Lie}}(Y/S,\mathcal M)\cdot Y
 =\underline{\operatorname{Lie}}(Y/S,\mathcal M)_Y,
\]
因而得到关于 \(\mathcal M\) 函子性的同构
\[
 L^u_{\underline{\operatorname{Hom}}_S(X,Y)/S}(\mathcal M)
 \xrightarrow{\sim}
 \underline{\operatorname{Hom}}_S
   (X,\underline{\operatorname{Lie}}(Y/S,\mathcal M)).
\]

对 \(S'\to S\), 以 \(u':X'\to Y'\) 表示 \(u\) 的基变换.
\(S\)-函子\footnote{编者注: 这是取 \(Z=Y\), 对 \(Y\)-群
\(u:X\to Y\) 和 \(T_{Y/S}(\mathcal M)\to Y\)
应用\SGARefLink{sga3:ii:definition:4-2-0-1}{定义 4.2.0.1}.}
\[
 \underline{\operatorname{Hom}}_{(Y/S)\text{-gr.}}
   \bigl(X,\underline{\operatorname{Lie}}(Y/S,\mathcal M)\cdot Y\bigr)
\]
由下式定义:
\[
\begin{aligned}
 &\underline{\operatorname{Hom}}_{(Y/S)\text{-gr.}}
 \bigl(X,\underline{\operatorname{Lie}}(Y/S,\mathcal M)\cdot Y\bigr)(S')\\
 &\quad=
 \operatorname{Hom}_{Y'\text{-gr.}}
 \bigl(X',(\underline{\operatorname{Lie}}(Y/S,\mathcal M)\cdot Y)_{S'}\bigr)\\
 &\quad=
 \operatorname{Hom}_{Y'\text{-gr.}}
 \bigl(X',\underline{\operatorname{Lie}}(Y'/S',\mathcal M)\cdot Y'\bigr).
\end{aligned}
\]
所以 \((\dagger)\) 诱导同构
\begin{equation}
 L^u_{\underline{\operatorname{Hom}}_{S\text{-gr.}}(X,Y)/S}
   (\mathcal M)
 \xrightarrow{\sim}
 \underline{\operatorname{Hom}}_{(Y/S)\text{-gr.}}
 \bigl(X,\underline{\operatorname{Lie}}(Y/S,\mathcal M)\cdot Y\bigr).
 \tag{\(\dagger'\)}\label{sga3:ii:subsection:4-2:equation:dagger-prime}
\end{equation}

另一方面, 复合态射

\SGAThreeNativeDiagram[0.50\textwidth]
  {figures/exp2/exp2_p027_adjoint_action_chain.tex}
  {\(X\) 在 \(\underline{\operatorname{Lie}}(Y/S,\mathcal M)\) 上的作用.}
  {fig:exp2-adjoint-action-chain}
  {Polo--Gille Exp2-14oct24.pdf, 组件内第 27 页; 合卷第 85 页; 印刷页码 75}

定义了 \(X\) 在 \(L=\underline{\operatorname{Lie}}(Y/S,\mathcal M)\) 上
通过群自同构给出的作用.

% Source page: Exp2 component-local p. 28; combined-reader p. 86; printed p. 76.
设 \(\Phi\in\underline{\operatorname{Hom}}_{Y/S}
(X,\underline{\operatorname{Lie}}(Y/S,\mathcal M)\cdot Y)\).
对 \(S''\to S'\to S\) 及 \(x\in X(S'')\), 存在唯一的表达式
\[
 \Phi(S')(x)=\varphi(S')(x)\cdot u'(x),
 \qquad
 \varphi(S')(x)\in
 \underline{\operatorname{Lie}}(Y'/S',\mathcal M)(S'').
\]
这确定了一个元素 \(\varphi\in\underline{\operatorname{Hom}}_S
(X,\underline{\operatorname{Lie}}(Y/S,\mathcal M))\).
映射 \(\Phi(S')\) 为群态射, 当且仅当对所有 \(x_1,x_2\in X(S'')\), 有
\[
\begin{aligned}
 \varphi(S')(x_1x_2)
 &=\varphi(S')(x_1)
   \bigl(u(x_1)\varphi(S')(x_2)u(x_1)^{-1}\bigr)\\
 &=\varphi(S')(x_1)
   \bigl(x_1\cdot\varphi(S')(x_2)\bigr),
\end{aligned}
\]
也就是说, 恰在 \(\varphi\in\underline Z^1_S
(X,\underline{\operatorname{Lie}}(Y/S,\mathcal M))\) 时成立.

\subsection*{命题 4.2.}
设 \(u:X\to Y\) 是 \(S\)-群态射. 存在关于 \(\mathcal M\) 函子性的
\(S\)-函子同构
\[
 L^u_{\underline{\operatorname{Hom}}_{S\text{-gr.}}(X,Y)/S}
   (\mathcal M)
 \xrightarrow{\sim}
 \underline Z^1_S
   \bigl(X,\underline{\operatorname{Lie}}(Y/S,\mathcal M)\bigr).
\]

再设 \(Y/S\) 满足 \(E\). 由 \SGARefLink{sga3:ii:corollary:4-2-0-3}{4.2.0.3},
与 \SGARefLink{sga3:ii:corollary:3-11-1}{3.11.1} 的证明一样, 可得
\(\underline{\operatorname{Hom}}_{S\text{-gr.}}(X,Y)/S\) 满足 \(E\).
因此 (参见 \SGARefLink{sga3:ii:subsection:3-5-1}{3.5.1}),
\[
\begin{aligned}
 L^u_{\underline{\operatorname{Hom}}_{S\text{-gr.}}(X,Y)/S}
 (\mathcal M\oplus\mathcal N)
 \simeq{}&
 L^u_{\underline{\operatorname{Hom}}_{S\text{-gr.}}(X,Y)/S}
 (\mathcal M)\\
 &\times_S
 L^u_{\underline{\operatorname{Hom}}_{S\text{-gr.}}(X,Y)/S}
 (\mathcal N).
\end{aligned}
\]
这也直接由\SGARefLink{sga3:ii:subsection:4-2}{命题~4.2} 及
\SGARefLink{sga3:ii:remark:4-2-b}{注~4.2.B}\(a\) 得到.
因此, \SGARefLink{sga3:ii:subsection:4-2}{命题~4.2} 中同构的两端都带有
由关于 \(\mathcal M\) 的函子性得到的 \(\mathcal O_S\)-模结构,
且该同构保持这些结构.

\SGAInlineRefTarget{sga3:ii:proposition:4-2-bis}\subsection*{\SGARefLink{sga3:ii:proposition:4-2-bis}{命题 4.2 bis}.}
设 \(u:X\to Y\) 是 \(S\)-群态射, 且 \(Y/S\) 满足 \(E\).
存在关于 \(\mathcal M\) 函子性的 \(\mathcal O_S\)-模同构
\[
 L^u_{\underline{\operatorname{Hom}}_{S\text{-gr.}}(X,Y)/S}
   (\mathcal M)
 \xrightarrow{\sim}
 \underline Z^1_S
   \bigl(X,\underline{\operatorname{Lie}}(Y/S,\mathcal M)\bigr).
\]
\footnote{编者注: 从此处起对原文作了扩充; 特别是增添了命题 4.2 bis.}

当 \(Y/S\) 满足 \(E\) 时, \SGARefLink{sga3:ii:corollary:3-13-1}{3.13.1}
也给出如下同构, 正如 \SGARefLink{sga3:ii:corollary:3-13-2}{3.13.2} 的证明中那样:
对每个 \(u\in\underline{\operatorname{Isom}}_{S\text{-gr.}}(X,Y)\),
存在关于 \(\mathcal M\) 函子性的同构
\begin{equation}
 L^u_{\underline{\operatorname{Isom}}_{S\text{-gr.}}(X,Y)/S}
   (\mathcal M)
 \xrightarrow{\sim}
 L^u_{\underline{\operatorname{Hom}}_{S\text{-gr.}}(X,Y)/S}
   (\mathcal M).
 \tag{*}\label{sga3:ii:subsection:4-2:equation:star}
\end{equation}
由此得到以下两个推论.

\SGAInlineRefTarget{sga3:ii:corollary:4-2-1}\subsection*{\SGARefLink{sga3:ii:corollary:4-2-1}{推论 4.2.1}.}
设 \(u:X\to Y\) 是 \(S\)-群态射. 若 \(Y/S\) 满足 \(E\),
则存在关于 \(\mathcal M\) 函子性的 \(\mathcal O_S\)-模同构
\[
 L^u_{\underline{\operatorname{Isom}}_{S\text{-gr.}}(X,Y)/S}
   (\mathcal M)
 \xrightarrow{\sim}
 \underline Z^1_S
   \bigl(X,\underline{\operatorname{Lie}}(Y/S,\mathcal M)\bigr).
\]

% Source page: Exp2 component-local p. 29; combined-reader p. 87; printed p. 77.
\SGAInlineRefTarget{sga3:ii:corollary:4-2-2}\subsection*{\SGARefLink{sga3:ii:corollary:4-2-2}{推论 4.2.2}.}
设 \(X\) 是 \(S\)-群. 若 \(X/S\) 满足 \(E\), 则存在关于 \(\mathcal M\)
函子性的 \(\mathcal O_S\)-模同构
\[
 \underline{\operatorname{Lie}}
   \bigl(\underline{\operatorname{Aut}}_{S\text{-gr.}}(X)/S,
         \mathcal M\bigr)
 \xrightarrow{\sim}
 \underline Z^1_S
   \bigl(X,\underline{\operatorname{Lie}}(X/S,\mathcal M)\bigr).
\]

若 \(Y\) 交换, 则其在 \(L=\underline{\operatorname{Lie}}(Y/S,\mathcal M)\)
上的伴随作用是平凡的, 所以
\(\underline Z^1_S(X,L)=
\underline{\operatorname{Hom}}_{S\text{-gr.}}(X,L)\).
\footnote{编者注: 以下这句话为增补内容.}

\SGAInlineRefTarget{sga3:ii:corollary:4-2-3}\subsection*{\SGARefLink{sga3:ii:corollary:4-2-3}{推论 4.2.3}.}
设 \(Y\) 是交换 \(S\)-群. 存在关于 \(\mathcal M\) 函子性的
\(S\)-函子同构
\[
 L^u_{\underline{\operatorname{Hom}}_{S\text{-gr.}}(X,Y)/S}
   (\mathcal M)
 \xrightarrow{\sim}
 \underline{\operatorname{Hom}}_{S\text{-gr.}}
   \bigl(X,\underline{\operatorname{Lie}}(Y/S,\mathcal M)\bigr).
\]

\subsection*{\SGAInlineRefTarget{sga3:ii:subsection:4-3}4.3.}
现在考虑 \(X\) 与 \(Y\) 是 \(\mathcal O_S\)-模的情形.
回顾 (参见 \SGARefLink{sga3:ii:definition:4-2-0-4}{4.2.0.4}),
\(T'_{Y/S}(\mathcal M)\), 相应地 \(\underline{\operatorname{Lie}}'(Y/S,\mathcal M)\),
表示函子 \(T_{Y/S}(\mathcal M)\), 相应地
\(\underline{\operatorname{Lie}}(Y/S,\mathcal M)\),
并在其上配备由 \(Y\) 的模结构诱导的 \(\mathcal O_S\)-模结构.

当 \(Y/S\) 满足 \(E\) 时, 我们总用
\(\underline{\operatorname{Lie}}(Y/S,\mathcal M)\) 表示该函子,
而给它配备的是对每个满足 \(E\) 的函子都已定义的 \(\mathcal O_S\)-模结构.
由 \SGARefLink{sga3:ii:proposition:3-9}{3.9},
\(\underline{\operatorname{Lie}}(Y/S,\mathcal M)\) 与
\(\underline{\operatorname{Lie}}'(Y/S,\mathcal M)\) 的底交换群结构相同,
但其模结构未必相同 (见 \SGARefLink{sga3:ii:subsection:6-3}{6.3} 的反例).
对 \(S'\to S\), \(a\in\mathcal O(S')\) 和
\(m\in\underline{\operatorname{Lie}}(Y/S,\mathcal M)(S')\),
用 \(a\mathbin{\cdot'}m\), 相应地 \(a\cdot m\),
表示来自 \(Y\) 的作用, 相应地表示由条件 \(E\) 得到的作用.
对 \(T_{Y/S}(\mathcal M)\) 上的这两种作用, 也采用同一约定.

存在 \(\mathcal O_S\)-模同构
\[
 T'_{Y/S}(\mathcal M)
 \simeq \underline{\operatorname{Lie}}'(Y/S,\mathcal M)\oplus Y.
\]
因此, 与\SGARefLink{sga3:ii:subsection:4-2}{命题~4.2} 和
\SGARefLink{sga3:ii:proposition:4-2-bis}{4.2 bis} 完全一样, 得到:

\subsection*{命题 4.3.}
设 \(u:X\to Y\) 是 \(\mathcal O_S\)-模态射.
存在关于 \(\mathcal M\) 函子性的 \(S\)-函子同构
\begin{equation}
 L^u_{\underline{\operatorname{Hom}}_{\mathcal O_S\text{-mod.}}(X,Y)/S}
   (\mathcal M)
 \xrightarrow{\sim}
 \underline{\operatorname{Hom}}_{\mathcal O_S\text{-mod.}}
   \bigl(X,\underline{\operatorname{Lie}}'(Y/S,\mathcal M)\bigr).
 \tag{*}\label{sga3:ii:subsection:4-3:equation:hom-lie-isomorphism-star}
\end{equation}
若 \(Y/S\) 满足 \(E\), 则
\(\underline{\operatorname{Hom}}_{\mathcal O_S\text{-mod.}}(X,Y)/S\) 满足 \(E\),
而当两端配备由关于 \(\mathcal M\) 的函子性得到的结构时,
\((*)\) 是 \(\mathcal O_S\)-模同构.
\footnote{编者注: 以下这句话为增补内容.}
\footnote{编者注: 右端的标量作用为 \((a\varphi)(x)=a\varphi(x)\),
其中 \(a\varphi(x)\) 表示 \(a\) 在
\(\varphi(x)\in
\underline{\operatorname{Lie}}(Y/S,\mathcal M)\) 上的作用.
这不同于作用 \((a\mathbin{\cdot'}\varphi)(x)
=a\mathbin{\cdot'}\varphi(x)=\varphi(ax)\),
但当 \(\underline{\operatorname{Lie}}(Y/S,\mathcal M)
=\underline{\operatorname{Lie}}'(Y/S,\mathcal M)\) 时, 两者相同.}

\SGAInlineRefTarget{sga3:ii:remark:4-3-bis}\subsection*{\SGARefLink{sga3:ii:remark:4-3-bis}{注 4.3 bis}.}
设 \(u:X\to Y\) 是 \(\mathcal O_S\)-模态射.
\footnote{编者注: 此注为增补内容; 它将在 \SGARefLink{sga3:ii:scholium:4-5-1}{4.5.1} 中用到.}
以 \(\tau_u\) 表示将每个 \(\mathcal O_S\)-线性态射
\(\varphi:X\to\underline{\operatorname{Lie}}'(Y/S,\mathcal M)\) 对应到
\[
 u\oplus\varphi:X\longrightarrow
 T'_{Y/S}(\mathcal M)
 =Y\oplus\underline{\operatorname{Lie}}'(Y/S,\mathcal M).
 \]
的映射.

% Source page: Exp2 component-local p. 30; combined-reader p. 88; printed p. 78.
\SGARefLink{sga3:ii:subsection:4-3}{4.3} 的同构包含在以下关于 \(\mathcal M\)
函子性的交换图中.

\SGAThreeNativeDiagram[0.70\textwidth]
  {figures/exp2/exp2_p030_tau_u_square.tex}
  {\SGARefLink{sga3:ii:subsection:4-3}{4.3} 的同构与 \(\tau_u\) 的相容性.}
  {fig:exp2-tau-u-square}
  {Polo--Gille Exp2-14oct24.pdf, 组件内第 30 页; 合卷第 88 页; 印刷页码 78}

此外, 当 \(Y/S\) 满足 \(E\) 时, 与
\SGARefLink{sga3:ii:corollary:3-13-2}{3.13.2} 的证明一样,
\SGARefLink{sga3:ii:corollary:3-13-1}{3.13.1} 给出:
对每个 \(u\in\underline{\operatorname{Isom}}_{\mathcal O_S\text{-mod.}}(X,Y)\), 有
\begin{equation}
 L^u_{\underline{\operatorname{Isom}}_{\mathcal O_S\text{-mod.}}(X,Y)/S}
   (\mathcal M)
 =
 L^u_{\underline{\operatorname{Hom}}_{\mathcal O_S\text{-mod.}}(X,Y)/S}
   (\mathcal M).
 \tag{*}\label{sga3:ii:subsection:4-3:equation:isom-hom-lie-equality-star}
\end{equation}

\SGAInlineRefTarget{sga3:ii:corollary:4-3-1}\subsection*{\SGARefLink{sga3:ii:corollary:4-3-1}{推论 4.3.1}.}
设 \(X\) 是 \(S\) 上满足 \(E\) 的 \(\mathcal O_S\)-模.
存在关于 \(\mathcal M\) 函子性的同构
\[
 \underline{\operatorname{Lie}}
  \bigl(\underline{\operatorname{Aut}}_{\mathcal O_S\text{-mod.}}(X)/S,
        \mathcal M\bigr)
 \xrightarrow{\sim}
 \underline{\operatorname{Hom}}_{\mathcal O_S\text{-mod.}}
  \bigl(X,\underline{\operatorname{Lie}}'(X/S,\mathcal M)\bigr),
\]
且这一同构保持由关于 \(\mathcal M\) 的函子性得到的 \(\mathcal O_S\)-模结构.
特别地, \(\underline{\operatorname{Aut}}_{\mathcal O_S\text{-mod.}}(X)/S\)
满足 \(E\).
\footnote{编者注: 比较编者注 (61). 最后一句话及其证明也是增补内容.}

\paragraph{证明.}
第一个断言由 \((*)\) 和 \SGARefLink{sga3:ii:subsection:4-3}{4.3} 得到.
对于第二个断言, 由于 \(X/S\) 满足 \(E\), 存在 \(\mathcal O_S\)-模同构
\[
 \underline{\operatorname{Lie}}'(X/S,\mathcal M\oplus\mathcal N)
 \simeq
 \underline{\operatorname{Lie}}'(X/S,\mathcal M)
 \times_S
 \underline{\operatorname{Lie}}'(X/S,\mathcal N).
\]
因此,
\[
\begin{aligned}
 &\underline{\operatorname{Lie}}
  \bigl(\underline{\operatorname{Aut}}_{\mathcal O_S\text{-mod.}}(X)/S,
  \mathcal M\oplus\mathcal N\bigr)\\
 &\quad\simeq
 \underline{\operatorname{Lie}}
  \bigl(\underline{\operatorname{Aut}}_{\mathcal O_S\text{-mod.}}(X)/S,
  \mathcal M\bigr)
 \times_S
 \underline{\operatorname{Lie}}
  \bigl(\underline{\operatorname{Aut}}_{\mathcal O_S\text{-mod.}}(X)/S,
  \mathcal N\bigr).
\end{aligned}
\]
现在由\SGARefLink{sga3:ii:remark:4-1-1-2}{注~4.1.1.2}\(a\) 即得断言.

\subsection*{\SGAInlineRefTarget{sga3:ii:subsection:4-3-2}4.3.2.}
沿这一方向继续讨论之前, 先更仔细地考察 \(Y\),
\(\underline{\operatorname{Lie}}(Y/S)\) 和
\(\underline{\operatorname{Lie}}'(Y/S)\) 之间的关系. 首先,
\begin{equation}
 \underline{\operatorname{Lie}}(\mathcal O_S/S,\mathcal M)
 =\underline{\operatorname{Lie}}'(\mathcal O_S/S,\mathcal M)
 =\mathbf W(\mathcal M),
 \tag{1}\label{sga3:ii:subsection:4-3-2:equation:1}
\end{equation}
其中 \(\mathbf W(\mathcal M)\) 的定义见 \SGARefLink{sga3:i:subsection:4-6}{I~4.6}.
因此有典范同构
\begin{equation}
 d:\mathcal O_S\xrightarrow{\sim}
 \underline{\operatorname{Lie}}(\mathcal O_S/S).
 \tag{2}\label{sga3:ii:subsection:4-3-2:equation:2}
\end{equation}

设 \(F\) 是 \(\mathcal O_S\)-模. 对每个 \(S_2\to S_1\to S\), 有一个双同态
\begin{equation}
 \left\{
 \begin{aligned}
   F(S_1)&\longrightarrow F(S_2),\\
   \mathcal O(S_1)&\longrightarrow\mathcal O(S_2),
 \end{aligned}
 \right.
 \tag{3}\label{sga3:ii:subsection:4-3-2:equation:3}
\end{equation}
因而有 \(\mathcal O(S_2)\)-模态射
\[
 F(S_1)\otimes_{\mathcal O(S_1)}\mathcal O(S_2)
 \longrightarrow F(S_2).
\]
\footnote{编者注: 这里将 \(S'\to S\) 换成 \(S_2\to S_1\to S\),
因为下文需要同时改变 \(S_2\) 和 \(S_1\). 下文还对原文作了扩充.}

% Source page: Exp2 component-local p. 31; combined-reader p. 89; printed p. 79.
特别地, 取 \(S_1=S'\) 与 \(S_2=I_{S'}(\mathcal M)\).
得到关于 \(\mathcal M\) 函子性的 \(\mathcal O(S')\)-模态射
\[
 F(S')\otimes_{\mathcal O(S')}
 T_{\mathcal O_S/S}(\mathcal M)(S')
 \longrightarrow T'_{F/S}(\mathcal M)(S').
\]
令 \(S'\) 改变, 这些态射给出 \(\mathcal O_S\)-模态射
\begin{equation}
 F\otimes_{\mathcal O_S}T_{\mathcal O_S/S}(\mathcal M)
 \longrightarrow T'_{F/S}(\mathcal M),
 \tag{4}\label{sga3:ii:subsection:4-3-2:equation:4}
\end{equation}
它们关于 \(\mathcal M\) 是函子性的. 由与切丛到底的投影的相容性, 继而得到态射
\begin{equation}
 F\otimes_{\mathcal O_S}
 \underline{\operatorname{Lie}}(\mathcal O_S/S,\mathcal M)
 \longrightarrow
 \underline{\operatorname{Lie}}'(F/S,\mathcal M),
 \tag{5}\label{sga3:ii:subsection:4-3-2:equation:5}
\end{equation}
且下图交换.

\SGAThreeNativeDiagram[0.78\textwidth]
  {figures/exp2/exp2_p031_tangent_exact_rows.tex}
  {分裂切序列中的态射 (4) 与 (5).}
  {fig:exp2-tangent-exact-rows}
  {Polo--Gille Exp2-14oct24.pdf, 组件内第 31 页; 合卷第 89 页; 印刷页码 79}

态射 (5) 可视为交换 \(S\)-群态射
\begin{equation}
 F\otimes_{\mathcal O_S}
 \underline{\operatorname{Lie}}(\mathcal O_S/S,\mathcal M)
 \longrightarrow
 \underline{\operatorname{Lie}}(F/S,\mathcal M).
 \tag{6}\label{sga3:ii:subsection:4-3-2:equation:6}
\end{equation}
将 \(F\) 与同构 \(d:\mathcal O_S\xrightarrow{\sim}
\underline{\operatorname{Lie}}(\mathcal O_S/S)\) 作张量积,
当 \(\mathcal M=\mathcal O_S\) 时, 得到交换 \(S\)-群态射
\begin{equation}
 F\xrightarrow{\sim}
 F\otimes_{\mathcal O_S}
 \underline{\operatorname{Lie}}(\mathcal O_S/S)
 \longrightarrow \underline{\operatorname{Lie}}(F/S),
 \tag{7}\label{sga3:ii:subsection:4-3-2:equation:7}
\end{equation}
也记为 \(d:F\to\underline{\operatorname{Lie}}(F/S)\).

\SGAInlineRefTarget{sga3:ii:remark:4-3-3}\subsection*{\SGARefLink{sga3:ii:remark:4-3-3}{注 4.3.3}.}
当 \(F/S\) 满足 \(E\) 时, 如果两端均配备由条件 \(E\) 从 \(\mathcal M\)
得到的模结构, 则态射 \SGARefLink{sga3:ii:subsection:4-3-2:equation:6}{(6)} 和 (7)
未必是 \(\mathcal O_S\)-线性的.

例如, 设 \(k\) 是特征为 \(p>0\) 的域, \(S=\operatorname{Spec}(k)\),
并设 \(F\) 是如下 \(\mathcal O_S\)-模: 它把每个 \(S\)-概形 \(T\)
对应到集合 \(F(T)=\Gamma(T,\mathcal O_T)\), 标量通过 \(p\) 次幂作用:
\[
 r\cdot f=r^pf
 \qquad(r\in\mathcal O(T),\ f\in F(T)).
\]
作为取值于群的 \(S\)-函子, \(F\simeq\mathbf G_{a,S}\).
因此 \(F\) 满足 \(E\), 且
\(\underline{\operatorname{Lie}}(F/S)\simeq
\underline{\operatorname{Lie}}(\mathbf G_{a,S}/S)\simeq\mathcal O_S\).
对每个 \(T\to S\), 典范态射 \(d:F(T)\to\mathcal O(T)\) 是恒等映射:
它保持交换群结构, 但不保持 \(\mathcal O_S\)-模结构.
\footnote{编者注: 原文断言, 当 \(F/S\) 满足 \(E\) 时,
(6) 以及由此得到的 (7) 都是 \(\mathcal O_S\)-模态射.
\SGARefLink{sga3:ii:subsection:6-3}{6.3} 中的反例与此矛盾.
因此删去了这一断言, 插入了前面的例子, 并相应修改了定义 4.4.
这不影响后续内容.}

% Source page: Exp2 component-local p. 32; combined-reader p. 90; printed p. 80.
\SGAInlineRefTarget{sga3:ii:remark:4-3-4}\subsection*{\SGARefLink{sga3:ii:remark:4-3-4}{注 4.3.4}.}
态射 (4) 与 (5) 可明确写成如下形式. 态射
\[
 \Theta:F\otimes_{\mathcal O_S}T_{\mathcal O_S/S}(\mathcal M)
 \longrightarrow T'_{F/S}(\mathcal M)
 =\underline{\operatorname{Hom}}_S(I_S(\mathcal M),F)
\]
对 \(S'\to S\), \(\alpha\in\mathcal O(I_{S'}(\mathcal M))\)
和 \(f:S'\to F\), 由下式定义:
\[
 \Theta(f\otimes\alpha)
 =\alpha\cdot(\tau_0\circ f)
 =\alpha\cdot(f\circ\rho),
\]
其中 \(\tau_0:F\to T'_{F/S}(\mathcal M)\) 是零截面,
\(\rho:I_{S'}(\mathcal M)\to S'\) 是结构态射.

态射 \(\Theta\) 诱导
\[
 \theta:F\otimes_{\mathcal O_S}
 \underline{\operatorname{Lie}}(\mathcal O_S/S,\mathcal M)
 \longrightarrow
 \underline{\operatorname{Lie}}'(F/S,\mathcal M).
\]
这由 (4) 之后指出的关于 \(\mathcal M\) 的函子性得出, 也可以直接看出. 事实上,
\[
 \underline{\operatorname{Lie}}'(F/S,\mathcal M)(S')
 =\{\,\varphi\in\operatorname{Hom}_S(I_{S'}(\mathcal M),F)
       \mid \varphi\circ\varepsilon_{\mathcal M}=e\,\},
\]
其中 \(e:S'\to F\) 是单位截面. 另一方面,
\(\underline{\operatorname{Lie}}(\mathcal O_S/S,\mathcal M)(S')\)
等于 \(\Gamma(S',\mathcal M)\), 后者是增广映射
\[
 \eta:\mathcal O(I_{S'}(\mathcal M))\longrightarrow\mathcal O(S').
\]
的核. 还需说明, 若 \(f\in F(S')\) 且 \(\alpha\in\Gamma(S',\mathcal M)\),
则 \(\alpha\cdot(f\circ\rho)\circ\varepsilon_{\mathcal M}=e\).

对 \(S\)-态射 \(\varepsilon_{\mathcal M}:S'\to I_{S'}(\mathcal M)\)
应用双同态 (3), 下图交换.

\SGAThreeNativeDiagram[0.43\textwidth]
  {figures/exp2/exp2_p032_epsilon_action_square.tex}
  {标量作用与 \(F(\varepsilon_{\mathcal M})\) 的相容性.}
  {fig:exp2-epsilon-action-square}
  {Polo--Gille Exp2-14oct24.pdf, 组件内第 32 页; 合卷第 90 页; 印刷页码 80}

于是, 对每个 \(\varphi:I_{S'}(\mathcal M)\to F\),
\[
 (\alpha\cdot\varphi)\circ\varepsilon_{\mathcal M}
 =\eta(\alpha)\cdot(\varphi\circ\varepsilon_{\mathcal M}),
\]
因此当 \(\eta(\alpha)=0\) 时,
\((\alpha\cdot\varphi)\circ\varepsilon_{\mathcal M}=e\).

特别地, 取 \(\mathcal M=t\mathcal O_S\), 有
\(\underline{\operatorname{Lie}}(\mathcal O_S/S)=t\mathcal O_S\),
而态射 \(F\to\underline{\operatorname{Lie}}'(F/S)\) 为
\(f\mapsto t\cdot(f\circ\rho)\).

\footnote{编者注: 此注为增补内容; 它将在 \SGARefLink{sga3:ii:lemma:4-6-2}{4.6.2} 中用到.}

\SGAInlineRefTarget{sga3:ii:remark:4-3-5}\subsection*{\SGARefLink{sga3:ii:remark:4-3-5}{注 4.3.5}.}
设 \(F\) 是 \(\mathcal O_S\)-模, 置
\[
 E=\underline{\operatorname{End}}_{\mathcal O_S\text{-mod.}}(F),
\]
并以 \(d_F\) 和 \(d_E\) 表示
\SGARefLink{sga3:ii:subsection:4-3-2:equation:5}{4.3.2(5)} 给出的
\(\mathcal O_S\)-线性态射:
\[
\begin{aligned}
 d_F:
 F\otimes_{\mathcal O_S}
 \underline{\operatorname{Lie}}(\mathcal O_S/S,\mathcal M)
 &\longrightarrow
 \underline{\operatorname{Lie}}'(F/S,\mathcal M),\\
 d_E:
 E\otimes_{\mathcal O_S}
 \underline{\operatorname{Lie}}(\mathcal O_S/S,\mathcal M)
 &\longrightarrow
 \underline{\operatorname{Lie}}'(E/S,\mathcal M).
\end{aligned}
\]
\SGARefLink{sga3:ii:remark:4-3-4}{注~4.3.4} 给出下列由
\(\mathcal O_S\)-线性态射组成的交换图.

\SGAThreeNativeDiagram[0.76\textwidth]
  {figures/exp2/exp2_p032_differentials_square.tex}
  {\(d_E\) 与 \(d_F\) 的相容性.}
  {fig:exp2-differentials-square}
  {Polo--Gille Exp2-14oct24.pdf, 组件内第 32 页; 合卷第 90 页; 印刷页码 80}

右边的竖直箭头是 \SGARefLink{sga3:ii:subsection:4-3}{4.3} 的同构 \((*)\).
因此, 由该结果可知, 若 \(F/S\) 满足 \(E\), 则 \(E/S\) 满足 \(E\),
而当右端各项配备由 \(E\) 得到的模结构时, \((*)\) 也是 \(\mathcal O_S\)-模同构.
\footnote{编者注: 此注为增补内容; 它将在 \SGARefLink{sga3:ii:proposition:4-5}{4.5}
和 \SGARefLink{sga3:ii:theorem:4-7}{4.7} 中用到.}

% Source page: Exp2 component-local p. 33; combined-reader p. 91; printed p. 81.
\SGAInlineRefTarget{sga3:ii:remark:4-4-0}\subsection*{\SGARefLink{sga3:ii:remark:4-4-0}{注 4.4.0}.}
在 \SGARefLink{sga3:ii:subsection:4-3-2}{4.3.2} 中,
态射 (4) 为同构当且仅当态射 \SGARefLink{sga3:ii:subsection:4-3-2:equation:5}{(5)} 为同构.
若这些等价条件成立, 则 \(F/S\) 满足 \(E\). 事实上, 只需验证
\[
 \underline{\operatorname{Lie}}'(F/S,\mathcal M\oplus\mathcal N)
 \simeq
 \underline{\operatorname{Lie}}'(F/S,\mathcal M)
 \times_S
 \underline{\operatorname{Lie}}'(F/S,\mathcal N).
\]
在假设下, 以下交换图中的水平箭头都是同构.

\SGAThreeNativeDiagram[0.84\textwidth]
  {figures/exp2/exp2_p033_good_module_square.tex}
  {用于验证条件 \(E\) 的直和比较.}
  {fig:exp2-good-module-square}
  {Polo--Gille Exp2-14oct24.pdf, 组件内第 33 页; 合卷第 91 页; 印刷页码 81}

所以第二个竖直箭头也是同构, 这恰好就是 \(F/S\) 的条件 \(E\).
\footnote{编者注: 鉴于定义 4.4 所作的修改 (参见编者注 (65)),
此处插入了这条注; 它原先出现在 \SGARefLink{sga3:ii:corollary:4-4-1}{4.4.1} 的证明中.}

\SGAInlineRefTarget{sga3:ii:definition:4-4}\subsection*{\SGARefLink{sga3:ii:definition:4-4}{定义 4.4}.}
若态射
\[
 F\otimes_{\mathcal O_S}T_{\mathcal O_S/S}(\mathcal M)
 \longrightarrow T_{F/S}(\mathcal M),
\]
或等价地, 态射
\[
 F\otimes_{\mathcal O_S}
 \underline{\operatorname{Lie}}(\mathcal O_S/S,\mathcal M)
 \longrightarrow
 \underline{\operatorname{Lie}}(F/S,\mathcal M),
\]
是交换 \(S\)-群的同构 (从而 \(F/S\) 满足 (E)), 且它们还保持由条件 \(E\)
得到的 \(\mathcal O_S\)-模结构, 则称 \(\mathcal O_S\)-模 \(F\)
为\emph{良模 (good module)}.

\SGAInlineRefTarget{sga3:ii:corollary:4-4-1}\subsection*{\SGARefLink{sga3:ii:corollary:4-4-1}{推论 4.4.1}.}
设 \(F\) 是 \(\mathcal O_S\)-模, 考虑以下条件:
\footnote{编者注: 原文证明了 (i) 推出 (ii) 和 (iii);
此处增添了 (ii) 推出 (i) 的证明.}
\begin{enumerate}
  \item[(i)] \(F\) 是良 \(\mathcal O_S\)-模;
  \item[(ii)] \(F/S\) 满足 \(E\), 且
    \(d:F\to\underline{\operatorname{Lie}}(F/S)\) 是 \(\mathcal O_S\)-模同构;
  \item[(iii)]
    \(\underline{\operatorname{Lie}}(F/S,\mathcal M)
      =\underline{\operatorname{Lie}}'(F/S,\mathcal M)\).
\end{enumerate}
则
\[
  \textup{(i)}\Longleftrightarrow\textup{(ii)}
  \Longrightarrow\textup{(iii)}.
\]

\paragraph{证明.}
蕴含关系 \(\textup{(i)}\Rightarrow\textup{(ii)}\) 由定义得到.
为证明 \(\textup{(ii)}\Rightarrow\textup{(i)}\), 还需说明关于 \(\mathcal M\)
函子性的交换 \(S\)-群态射
\[
 F\otimes_{\mathcal O_S}
 \underline{\operatorname{Lie}}(\mathcal O_S/S,\mathcal M)
 \xrightarrow{\sim}
 \underline{\operatorname{Lie}}(F/S,\mathcal M)
\]
是 \(\mathcal O_S\)-模同构. 由于 \(F/S\) 满足 \(E\),
两端都将 \(\mathcal O_S\) 的有限个副本的直和变成交换 \(S\)-群的有限积.
因此, 归结到 \(\mathcal M=\mathcal O_S\) 的情形, 而这恰好就是假设.

最后, \(\textup{(i)}\Rightarrow\textup{(iii)}\) 由定义以及
\SGARefLink{sga3:ii:subsection:4-3-2:equation:5}{4.3.2(5)} 的同构
\[
 F\otimes_{\mathcal O_S}
 \underline{\operatorname{Lie}}(\mathcal O_S/S,\mathcal M)
 \xrightarrow{\sim}
 \underline{\operatorname{Lie}}'(F/S,\mathcal M)
\]
为 \(\mathcal O_S\)-线性这一事实得出.

\SGAInlineRefTarget{sga3:ii:examples:4-4-2}\SGAInlineRefTarget{sga3:ii:example:4-4-2}\subsection*{\SGARefLink{sga3:ii:example:4-4-2}{例 4.4.2}.}
对每个拟凝聚 \(\mathcal O_S\)-模 \(\mathcal E\),
\SGARefLink{sga3:i:subsection:4-6}{I~4.6} 中定义的 \(\mathcal O_S\)-模
\(\mathbf V(\mathcal E)\) 和 \(\mathbf W(\mathcal E)\) 都是良模.

% Source page: Exp2 component-local p. 34; combined-reader p. 92; printed p. 82.
事实上, 对每个 \(f:S'\to S\), 态射
\[
\begin{aligned}
 \mathbf V(\mathcal E)(S')\otimes_{\mathcal O(S')}
   \mathcal O(I_{S'}(\mathcal M))
 &\longrightarrow T_{\mathbf V(\mathcal E)/S}(\mathcal M)(S'),\\
 \mathbf W(\mathcal E)(S')\otimes_{\mathcal O(S')}
   \mathcal O(I_{S'}(\mathcal M))
 &\longrightarrow T_{\mathbf W(\mathcal E)/S}(\mathcal M)(S')
\end{aligned}
\]
分别对应于
\[
\begin{aligned}
 &\operatorname{Hom}_{\mathcal O_{S'}\text{-mod.}}
   \bigl(f^*(\mathcal E),\mathcal O_{S'}\bigr)
 \otimes_{\mathcal O(S')}
 \Gamma\bigl(S',D_{\mathcal O_{S'}}(\mathcal M)\bigr)\\
 &\hspace{5em}\longrightarrow
 \operatorname{Hom}_{\mathcal O_{S'}\text{-mod.}}
   \bigl(f^*(\mathcal E),D_{\mathcal O_{S'}}(\mathcal M)\bigr)
 \qquad\textup{(V)}
\\[-2pt]
 &\Gamma(S',f^*(\mathcal E))
 \otimes_{\mathcal O(S')}
 \Gamma\bigl(S',D_{\mathcal O_{S'}}(\mathcal M)\bigr)\\
 &\hspace{5em}\longrightarrow
 \Gamma\bigl(S',f^*(\mathcal E)
   \otimes_{\mathcal O_{S'}}D_{\mathcal O_{S'}}(\mathcal M)\bigr)
 \qquad\textup{(W)}.
\end{aligned}
\]
它们是同构, 因为作为 \(\mathcal O_{S'}\)-模,
\(D_{\mathcal O_{S'}}(\mathcal M)\) 同构于 \(\mathcal O_{S'}\) 的有限个副本的直和.
\footnote{编者注: 前面的证明是增补内容.}

% End of the assigned scope. Proposition 4.5 begins immediately here in
% Exp2-14oct24.pdf, component-local p. 34 (combined-reader p. 92).
